\honest{The World: An Introduction}{Rob Bensinger}{2015}

The earlier books compared humans, as goal-seeking systems, with evolution and artificial intelligence. This book compares humans, as physical systems, with processes that are not mind-like. We are built of inhuman parts. As Giulio Giorello put it in an interview with Daniel Dennett: ``Yes, we have a soul. But it's made of lots of tiny robots.''

I preview the book's sequences, and then say a little about two debates it enters: consciousness and quantum physics. Readers eager to skip ahead may.

Better models of bat behavior or bat neurology, I say, may not tell us what echolocation feels like to the bat. Could an unconscious automaton copy all the behavior of a conscious one? Philosophers call such automata ``zombies.'' \nb{The zombies of the argument against physicalism are copies in every physical respect, not only in behavior. The \textsc{Stanford Encyclopedia of Philosophy} notes that ``plenty of physicalists'' accept that a merely behavioral copy might lack experience. The introduction states the stronger version later.} Alien psychologists with a perfect model of our behavior might still not know what redness feels like to us.

Thomas Nagel and David Chalmers have argued that third-person models can never fully capture first-person consciousness. Chalmers holds that consciousness is a ``further fact'' beyond the physical facts. Many find this persuasive, so I ask: ``should we grant its conclusion and ditch physicalism?'' \nb{Nagel himself did not draw this conclusion: ``It would be a mistake to conclude that physicalism must be false.'' The argument against physicalism is Chalmers's.} We should not reject the argument because it sounds strange, and ``Physicalism 201'' will test it against a technical account of explanation and evidence.

Quantum mechanics has amplitudes that evolve deterministically under the Schrödinger equation, until someone looks and Born's probabilistic rule takes over. I quote John Bell asking whether the wavefunction had to wait for ``some better qualified system \ldots\ with a PhD.'' The Copenhagen school split: some said consciousness causes collapse, Heisenberg's followers said physics is about our knowledge, and a third tradition said ``shut up and calculate.'' \nb{The consciousness view is usually traced to von Neumann and Wigner; the \textsc{Stanford Encyclopedia} says such views ``were definitely not Bohr's.'' ``Shut up and calculate'' is David Mermin's summary of Copenhagen, from 1989.}

Yudkowsky uses this controversy to test ideas from the earlier sequences. Yudkowsky is not a physicist, and neither am I, so I list outside sources for checking the arguments. Many-worlds, I say, ``has become very popular in recent decades among physicists,'' though a number reject it or stay agnostic. \nb{No source is given. In the two small polls found, of conference participants in 1997 and 2011, many-worlds was popular but held by a minority.}

Last, a reassurance: most of physics can be understood without a settled view of the ultimate nature, and size, of the physical world.
